\scanpage{123}
\begin{center}- 92 -\end{center}
\margin{n\textsuperscript{o} 371}

\uline{de sous-modules}; \uline{pour tout} $i$, \uline{soit} $u_i=M_i\to N_i$ \uline{une application} $A$-\uline{linéaire et soit} $u=\bigoplus_{1\leq i\leq n}u_i:M\to N$. \uline{Alors le diagramme}
\diagram{\begin{tikzcd}[ampersand replacement=\&,column sep=large,row sep=large]
{}^{\mathrm g}\!\displaystyle\bigotimes_{1\leq i\leq n}(\bigwedge(M_i)) \arrow[r] \arrow[d,"{}^{\mathrm g}\!\bigotimes_{1\leq i\leq n}(\bigwedge(u_i))"'] \& \bigwedge(M) \arrow[d,"\bigwedge(u)"] \\
{}^{\mathrm g}\!\displaystyle\bigotimes_{1\leq i\leq n}(\bigwedge(N_i)) \arrow[r] \& \bigwedge(N)
\end{tikzcd}}
(\uline{où les flèches horizontales sont les isomorphismes définis dans la prop.7}) \uline{est commutatif}.

Il suffit de voir que les deux applications ${}^{\mathrm g}\!\bigotimes_i(\bigwedge(M_i))\to\bigwedge(N)$ définies par le diagramme prennent la même valeur dans chacun des modules $\bigotimes_{1\leq i\leq n}(\mathop{\bigwedge}\limits^{m_i}M_i)$, ce qui résulte aussitôt des définitions.

COROLLAIRE 3.- \uline{Soit} $u:M\to N$ \uline{une application} $A$-\uline{linéaire injective}, \uline{telle que} $u(M)$ \uline{soit facteur direct de} $N$. \uline{Alors l'homomorphisme} $\bigwedge(u):\bigwedge(M)\to\bigwedge(N)$ \uline{est injectif et} $\bigwedge(N)$ \uline{est somme directe de la sous-algèbre graduée} $\bigwedge(u)(\bigwedge(M))$ \uline{et d'un idéal bilatère gradué engendré par un supplémentaire de} $u(M)$ \uline{dans} $N$.

Cela résulte du cor.1 appliqué pour $n=2$, car la somme des sous-modules $(\mathop{\bigwedge}\limits^p M_1)\otimes(\mathop{\bigwedge}\limits^q M_2)$ pour tous les couples $(p,q)$ tels que $p\geq1$ n'est autre que l'idéal bilatère $m_1\otimes(\bigwedge(M_2))$ de $(\bigwedge(M_1))\otimes^{\mathrm g}(\bigwedge(M_2))$, où $m_1$ est la somme directe des $\bigwedge(M_1)$ pour $p\geq1$, idéal bilatère de $\bigwedge(M_1)$. On peut aussi donner une démonstration directe n'utilisant pas la prop.7 et calquée sur celle du § 4, n\textsuperscript{o} 2, prop.4.

COROLLAIRE 4.- \uline{Tout élément} $z$ \uline{de} $\bigwedge(M)$ \uline{dont la composante homogène scalaire} (i.e. \uline{de degré} 0) \uline{est nulle}, \uline{est nilpotent}.

Il est clair en effet que tout élément de $\bigwedge(M)$ appartient à une sous-algèbre
\margin{122}
